\chapter{INTRODUCCIÓN}
\label{cap:introduccion}

El proyecto aborda el diseño de una máquina para la limpieza de rodamientos de gran tamaño utilizados en equipos de rodillos de molienda de alta presión, conocidos como HPGR, en entornos mineros. Se sitúa en el mantenimiento mecánico, en una primera etapa de lavado exterior que prepara el rodamiento para la inspección visual y precede al desmontaje por componentes y a la inspección metrológica. Su desarrollo integra diseño de máquinas, resistencia de materiales, mecánica de fluidos, selección de materiales y seguridad, porque el problema combina el manejo de una carga pesada con la limpieza de superficies de precisión.

La evaluación del rodamiento incluye una observación inicial de las superficies accesibles para identificar daños manifiestos, como fracturas, abolladuras o rayaduras de gran tamaño. Esta primera revisión permite registrar la condición de recepción y orientar la inspección posterior. La grasa y los residuos pueden ocultar otros indicios de deterioro, por lo que la limpieza y el secado preceden al examen detallado de pistas, elementos rodantes y jaula. SKF distingue la revisión visual de las superficies expuestas de la inspección más detallada y, para los rodamientos retirados, indica limpiar, secar, inspeccionar y proteger contra la corrosión aquellos que pueden reutilizarse \cite{skfmaintenance2011}.

En aplicaciones HPGR, esta tarea se relaciona con condiciones exigentes de operación. Schaeffler documenta en una instalación minera peruana la presencia de cargas de choque y contaminación elevada, junto con exigencias específicas de lubricación de los rodamientos \cite{schaeffler2021}. Durante el mantenimiento, la remoción de grasa y partículas debe considerar el estado del componente y el acceso a sus partes. Un procedimiento adecuado necesita controlar la eliminación de los residuos y la protección posterior del rodamiento frente a una nueva contaminación.

El diseño toma como referencia el rodamiento SKF 241/900 ECAK30/W33 y adopta un diámetro interior de \SI{900}{\milli\metre}, un diámetro exterior de \SI[group-minimum-digits=4]{1420}{\milli\metre}, un ancho de \SI{515}{\milli\metre} y una masa aproximada de \SI[group-minimum-digits=4]{3032}{\kilogram} \cite{skf241}. Estas magnitudes delimitan el espacio de trabajo y la carga que debe recibir el soporte. También condicionan la ubicación del equipo de izaje, la transferencia del rodamiento y el acceso a las superficies que requieren limpieza.

El procedimiento que una empresa peruana de remanufactura de rodamientos aplica a un rodamiento HPGR destinado a ese fin consiste en desmontarlo y limpiar sus componentes a mano antes de la inspección metrológica y de los ensayos no destructivos \cite{taller_jefe_2026}. Según el jefe de taller, una campaña de cuatro rodamientos demanda unos diez días de trabajo en un turno con cuatro técnicos, consume cerca de 40 L de solvente y genera alrededor de 300 kg de residuos contaminados entre grasa y trapos \cite{taller_jefe_2026}. Por rodamiento, esto equivale a unos diez turnos-persona, 10 L de solvente y 75 kg de residuos. La limpieza de los elementos rodantes es la etapa más demandante, porque la grasa queda retenida entre los rodillos y los aros. La única lavadora industrial del taller no se usa con rodamientos HPGR porque los rodillos salían sucios \cite{taller_jefe_2026}. Durante la tarea, el técnico decide la distancia de aplicación, el tiempo en cada zona y cuándo da por limpia una superficie, con trajes desechables y guantes de nitrilo como única barrera frente al solvente y la grasa. Esta dependencia del operario impide fijar una secuencia repetible y comparar intervenciones. En el Perú operan alrededor de 30 molinos HPGR, con unos 120 rodamientos 241/900 instalados \cite{taller_ingeniero_2026}, y cada rodamiento que se retira para remanufactura pasa por este procedimiento. La necesidad práctica consiste en disponer de medios que permitan controlar esas variables y organizar la tarea con criterios verificables.

El alcance de esta máquina corresponde a una primera etapa de lavado exterior del rodamiento completo, retirado del HPGR pero aún montado, cuyo objetivo es dejarlo en condiciones para la inspección visual antes de su desmontaje por componentes y de la inspección metrológica y los ensayos no destructivos. La máquina trata las superficies exteriores y otras superficies accesibles definidas para esta etapa; no busca limpiar los espacios internos entre rodillos y aros ni sustituye la limpieza manual de componentes requerida para la remanufactura. La limpieza de rodillos, aros y jaula desmontados queda planteada como una segunda etapa conceptual. Sus útiles, trayectorias y criterios no se dimensionan ni evalúan en este proyecto, y la limpieza interna con el rodamiento montado queda fuera del alcance de la primera etapa.

El manejo de una pieza de más de tres toneladas exige definir cómo se recibe, se posiciona y se mantiene estable. SKF recomienda emplear medios de elevación y transporte adecuados para rodamientos grandes y pesados \cite{skfmaintenance2011}. La transferencia desde el equipo de izaje, el traslado hasta la posición de lavado y la inmovilización forman una misma secuencia, y cada etapa condiciona la exposición del personal a atrapamientos, golpes o movimientos imprevistos de la carga.

Los usuarios directos son los técnicos encargados de la limpieza y la inspección. Necesitan acceder al componente y operar la máquina sin depender de maniobras improvisadas. El área de planificación de mantenimiento requiere una secuencia definida para asignar personal, equipos y tiempo de intervención. Los responsables de seguridad y medio ambiente intervienen en el control de la exposición al fluido y en el manejo de los residuos. Para la gestión del taller también importan el costo de implementación, la disponibilidad de repuestos y la facilidad de mantenimiento de la propia máquina.

El entorno de taller impone condiciones sobre el proceso. El fluido que arrastra grasa y partículas debe recogerse para evitar su dispersión y facilitar su gestión. Su elección debe considerar los materiales del rodamiento, el tipo de contaminante y la protección superficial posterior. El ASM Handbook incluye estos factores, junto con el costo, las instalaciones disponibles y la manipulación de la pieza, entre los criterios para seleccionar un proceso de limpieza \cite{chalk1994}. Por ello, el sistema de aplicación necesita desarrollarse junto con la contención y el drenaje.

Estas condiciones generan decisiones de diseño que se afectan entre sí. Un soporte más amplio puede favorecer la estabilidad y ocupar espacio necesario para el acceso; una trayectoria con mayor cobertura puede exigir mecanismos adicionales y aumentar el costo. La asistencia al movimiento también debe evaluarse frente al esfuerzo requerido y a las condiciones de maniobra. El problema exige compatibilizar estructura, movimiento y circulación del fluido dentro de los recursos del taller. Esa interdependencia justifica comparar varias alternativas antes de definir la configuración de la máquina.

¿Cómo diseñar una máquina para limpiar rodamientos HPGR de gran tamaño en talleres mineros?

La utilidad técnica del proyecto consiste en organizar el soporte y la aplicación del fluido como funciones coordinadas, con condiciones de operación que puedan verificarse. Para el personal de mantenimiento, el valor esperado se encuentra en disponer de un procedimiento más consistente y reducir la intervención directa durante el lavado. La protección del rodamiento también orienta el diseño: los apoyos, los mecanismos de posicionamiento y el proceso de limpieza deben permitir su inspección sin introducir daños atribuibles al manejo del componente.

La propuesta aporta valor al taller cuando su operación es compatible con el espacio, los medios de izaje y los recursos disponibles. La contención del fluido y la recogida de residuos atienden las condiciones de trabajo y el impacto ambiental del proceso. La selección de componentes comerciales y el acceso para mantenimiento influyen en el costo de implementación y conservación. La evaluación económica estima la inversión completa de la máquina y el beneficio por rodamiento que haría falta para recuperarla.

\section*{Objetivo general}

Diseñar una máquina para la limpieza controlada de rodamientos HPGR de gran tamaño que integre soporte, posicionamiento, aplicación de fluido y contención en un taller de mantenimiento minero.

\section*{Objetivos específicos}

\begin{enumerate}
\item Definir los requisitos de diseño a partir de la geometría y masa del rodamiento, las necesidades de los usuarios y las restricciones del taller.
\item Generar alternativas que integren las funciones de soporte, traslado, aplicación del fluido y recogida de residuos.
\item Seleccionar la alternativa más adecuada mediante una matriz de decisión que considere desempeño, seguridad, costo y facilidad de fabricación y mantenimiento.
\item Dimensionar los subsistemas mecánicos y establecer los parámetros hidráulicos de la alternativa seleccionada mediante cálculos de ingeniería y datos técnicos de componentes.
\item Evaluar el desempeño calculado frente a los requisitos establecidos y reajustar la geometría, los componentes o las condiciones de operación cuando los resultados lo requieran.
\end{enumerate}

\section*{Alcance}

El proyecto llega hasta la ingeniería preliminar de la máquina, verificada por cálculo y simulación, con las actividades del cronograma del Anexo~\ref{anx:portafolio}. Se desarrollan con detalle la mesa de traslado y su bastidor, los apoyos y los centradores del rodamiento, el sistema de aspersión con su red hidráulica, la simulación del mojado, la seguridad, el presupuesto y el plan de pruebas por simulación. Los entregables son la comparación de alternativas, el modelo CAD del conjunto, los planos de conjunto y de los subconjuntos principales, el esquema hidráulico, las memorias de cálculo, el presupuesto de inversión y la matriz de cumplimiento de especificaciones.

Otros componentes necesarios para que la máquina funcione se resuelven con menor profundidad: el soporte superior del distribuidor, el motorreductor y el limitador de par, los anclajes y la revisión sísmica. Para ellos se eligieron productos de catálogo o se hicieron cálculos de comprobación, que se resumen en el Anexo~\ref{anx:calculos}. El calentamiento del baño, la motorización del traslado, el control eléctrico, el secado, el tratamiento de efluentes y la obra civil quedan como recomendaciones para la etapa siguiente. Las simulaciones valen para las cargas, la geometría y las hipótesis de cada modelo. Se excluyen la fabricación, la puesta en marcha y los ensayos físicos, incluida la comprobación experimental de la limpieza.

Las restricciones principales son la masa y las dimensiones del rodamiento de referencia; el espacio disponible para carga, traslado y lavado; la compatibilidad con los medios de izaje; la prevención de atrapamientos y exposición al fluido; la compatibilidad de los materiales con el agente de limpieza; la contención de residuos; y el costo y disponibilidad de los componentes. Estas condiciones delimitan la solución y justifican los compromisos entre desempeño, seguridad y viabilidad constructiva.
